% Options for packages loaded elsewhere
\PassOptionsToPackage{unicode}{hyperref}
\PassOptionsToPackage{hyphens}{url}
\documentclass[
]{article}
\usepackage{xcolor}
\usepackage[margin=1in]{geometry}
\usepackage{amsmath,amssymb}
\setcounter{secnumdepth}{-\maxdimen} % remove section numbering
\usepackage{iftex}
\ifPDFTeX
  \usepackage[T1]{fontenc}
  \usepackage[utf8]{inputenc}
  \usepackage{textcomp} % provide euro and other symbols
\else % if luatex or xetex
  \usepackage{unicode-math} % this also loads fontspec
  \defaultfontfeatures{Scale=MatchLowercase}
  \defaultfontfeatures[\rmfamily]{Ligatures=TeX,Scale=1}
\fi
\usepackage{lmodern}
\ifPDFTeX\else
  % xetex/luatex font selection
\fi
% Use upquote if available, for straight quotes in verbatim environments
\IfFileExists{upquote.sty}{\usepackage{upquote}}{}
\IfFileExists{microtype.sty}{% use microtype if available
  \usepackage[]{microtype}
  \UseMicrotypeSet[protrusion]{basicmath} % disable protrusion for tt fonts
}{}
\makeatletter
\@ifundefined{KOMAClassName}{% if non-KOMA class
  \IfFileExists{parskip.sty}{%
    \usepackage{parskip}
  }{% else
    \setlength{\parindent}{0pt}
    \setlength{\parskip}{6pt plus 2pt minus 1pt}}
}{% if KOMA class
  \KOMAoptions{parskip=half}}
\makeatother
\usepackage{color}
\usepackage{fancyvrb}
\newcommand{\VerbBar}{|}
\newcommand{\VERB}{\Verb[commandchars=\\\{\}]}
\DefineVerbatimEnvironment{Highlighting}{Verbatim}{commandchars=\\\{\}}
% Add ',fontsize=\small' for more characters per line
\usepackage{framed}
\definecolor{shadecolor}{RGB}{248,248,248}
\newenvironment{Shaded}{\begin{snugshade}}{\end{snugshade}}
\newcommand{\AlertTok}[1]{\textcolor[rgb]{0.94,0.16,0.16}{#1}}
\newcommand{\AnnotationTok}[1]{\textcolor[rgb]{0.56,0.35,0.01}{\textbf{\textit{#1}}}}
\newcommand{\AttributeTok}[1]{\textcolor[rgb]{0.13,0.29,0.53}{#1}}
\newcommand{\BaseNTok}[1]{\textcolor[rgb]{0.00,0.00,0.81}{#1}}
\newcommand{\BuiltInTok}[1]{#1}
\newcommand{\CharTok}[1]{\textcolor[rgb]{0.31,0.60,0.02}{#1}}
\newcommand{\CommentTok}[1]{\textcolor[rgb]{0.56,0.35,0.01}{\textit{#1}}}
\newcommand{\CommentVarTok}[1]{\textcolor[rgb]{0.56,0.35,0.01}{\textbf{\textit{#1}}}}
\newcommand{\ConstantTok}[1]{\textcolor[rgb]{0.56,0.35,0.01}{#1}}
\newcommand{\ControlFlowTok}[1]{\textcolor[rgb]{0.13,0.29,0.53}{\textbf{#1}}}
\newcommand{\DataTypeTok}[1]{\textcolor[rgb]{0.13,0.29,0.53}{#1}}
\newcommand{\DecValTok}[1]{\textcolor[rgb]{0.00,0.00,0.81}{#1}}
\newcommand{\DocumentationTok}[1]{\textcolor[rgb]{0.56,0.35,0.01}{\textbf{\textit{#1}}}}
\newcommand{\ErrorTok}[1]{\textcolor[rgb]{0.64,0.00,0.00}{\textbf{#1}}}
\newcommand{\ExtensionTok}[1]{#1}
\newcommand{\FloatTok}[1]{\textcolor[rgb]{0.00,0.00,0.81}{#1}}
\newcommand{\FunctionTok}[1]{\textcolor[rgb]{0.13,0.29,0.53}{\textbf{#1}}}
\newcommand{\ImportTok}[1]{#1}
\newcommand{\InformationTok}[1]{\textcolor[rgb]{0.56,0.35,0.01}{\textbf{\textit{#1}}}}
\newcommand{\KeywordTok}[1]{\textcolor[rgb]{0.13,0.29,0.53}{\textbf{#1}}}
\newcommand{\NormalTok}[1]{#1}
\newcommand{\OperatorTok}[1]{\textcolor[rgb]{0.81,0.36,0.00}{\textbf{#1}}}
\newcommand{\OtherTok}[1]{\textcolor[rgb]{0.56,0.35,0.01}{#1}}
\newcommand{\PreprocessorTok}[1]{\textcolor[rgb]{0.56,0.35,0.01}{\textit{#1}}}
\newcommand{\RegionMarkerTok}[1]{#1}
\newcommand{\SpecialCharTok}[1]{\textcolor[rgb]{0.81,0.36,0.00}{\textbf{#1}}}
\newcommand{\SpecialStringTok}[1]{\textcolor[rgb]{0.31,0.60,0.02}{#1}}
\newcommand{\StringTok}[1]{\textcolor[rgb]{0.31,0.60,0.02}{#1}}
\newcommand{\VariableTok}[1]{\textcolor[rgb]{0.00,0.00,0.00}{#1}}
\newcommand{\VerbatimStringTok}[1]{\textcolor[rgb]{0.31,0.60,0.02}{#1}}
\newcommand{\WarningTok}[1]{\textcolor[rgb]{0.56,0.35,0.01}{\textbf{\textit{#1}}}}
\usepackage{longtable,booktabs,array}
\usepackage{calc} % for calculating minipage widths
% Correct order of tables after \paragraph or \subparagraph
\usepackage{etoolbox}
\makeatletter
\patchcmd\longtable{\par}{\if@noskipsec\mbox{}\fi\par}{}{}
\makeatother
% Allow footnotes in longtable head/foot
\IfFileExists{footnotehyper.sty}{\usepackage{footnotehyper}}{\usepackage{footnote}}
\makesavenoteenv{longtable}
\usepackage{graphicx}
\makeatletter
\newsavebox\pandoc@box
\newcommand*\pandocbounded[1]{% scales image to fit in text height/width
  \sbox\pandoc@box{#1}%
  \Gscale@div\@tempa{\textheight}{\dimexpr\ht\pandoc@box+\dp\pandoc@box\relax}%
  \Gscale@div\@tempb{\linewidth}{\wd\pandoc@box}%
  \ifdim\@tempb\p@<\@tempa\p@\let\@tempa\@tempb\fi% select the smaller of both
  \ifdim\@tempa\p@<\p@\scalebox{\@tempa}{\usebox\pandoc@box}%
  \else\usebox{\pandoc@box}%
  \fi%
}
% Set default figure placement to htbp
\def\fps@figure{htbp}
\makeatother
\setlength{\emergencystretch}{3em} % prevent overfull lines
\providecommand{\tightlist}{%
  \setlength{\itemsep}{0pt}\setlength{\parskip}{0pt}}
\usepackage{booktabs}
\usepackage{longtable}
\usepackage{array}
\usepackage{multirow}
\usepackage{wrapfig}
\usepackage{float}
\usepackage{colortbl}
\usepackage{pdflscape}
\usepackage{tabu}
\usepackage{threeparttable}
\usepackage{threeparttablex}
\usepackage[normalem]{ulem}
\usepackage{makecell}
\usepackage{xcolor}
\usepackage{bookmark}
\IfFileExists{xurl.sty}{\usepackage{xurl}}{} % add URL line breaks if available
\urlstyle{same}
\hypersetup{
  pdftitle={10.RMarkdown},
  pdfauthor={Carter Stancil},
  hidelinks,
  pdfcreator={LaTeX via pandoc}}

\title{10.RMarkdown}
\author{Carter Stancil}
\date{2026-04-10}

\begin{document}
\maketitle

\subsection{R Markdown}\label{r-markdown}

This is an R Markdown document. Markdown is a simple formatting syntax
for authoring HTML, PDF, and MS Word documents. For more details on
using R Markdown see \url{http://rmarkdown.rstudio.com}.

When you click the \textbf{Knit} button a document will be generated
that includes both content as well as the output of any embedded R code
chunks within the document. You can embed an R code chunk like this:

\subsection{Including Plots}\label{including-plots}

You can also embed plots, for example:

\begin{Shaded}
\begin{Highlighting}[]
\CommentTok{\# This option sets the R output in your document to not have \#\# before it, which}
\CommentTok{\#  is the default. You can customise pretty much anything (see the R Markdown}
\CommentTok{\#  book).}
\NormalTok{knitr}\SpecialCharTok{::}\NormalTok{opts\_chunk}\SpecialCharTok{$}\FunctionTok{set}\NormalTok{(}\AttributeTok{comment =} \StringTok{""}\NormalTok{)}
\CommentTok{\# \textless{}!{-}{-} Note that because the above chunk is R code, the \textasciigrave{}\# This option...\textasciigrave{} represents a comment because of the \textasciigrave{}\#\textasciigrave{}. Your text editor should hopefully colour code this file so you can distinguish R Markdown headings and R comments. And, yes, "chunk" is the technical name for these chunks of code. {-}{-}\textgreater{}}

\DocumentationTok{\#\# Generate data}

\CommentTok{\# First we\textquotesingle{}ll need some libraries:}
\end{Highlighting}
\end{Shaded}

\section{If you do not have any of them then install them by typing, for
example,}\label{if-you-do-not-have-any-of-them-then-install-them-by-typing-for-example}

\begin{Shaded}
\begin{Highlighting}[]
\FunctionTok{install.packages}\NormalTok{(kableExtra)}
\end{Highlighting}
\end{Shaded}

in your R console (or install using RStudio if you use that).

\begin{Shaded}
\begin{Highlighting}[]
\FunctionTok{library}\NormalTok{(kableExtra)}
\end{Highlighting}
\end{Shaded}

\begin{verbatim}
Warning: package 'kableExtra' was built under R version 4.5.2
\end{verbatim}

\begin{Shaded}
\begin{Highlighting}[]
\FunctionTok{library}\NormalTok{(tibble)}
\FunctionTok{library}\NormalTok{(xtable)}
\end{Highlighting}
\end{Shaded}

\begin{verbatim}
Warning: package 'xtable' was built under R version 4.5.2
\end{verbatim}

Now generate some data:

\begin{Shaded}
\begin{Highlighting}[]
\FunctionTok{set.seed}\NormalTok{(}\DecValTok{42}\NormalTok{)}
\NormalTok{n }\OtherTok{\textless{}{-}} \DecValTok{30}                      \CommentTok{\# sample size}
\NormalTok{x }\OtherTok{\textless{}{-}} \DecValTok{1}\SpecialCharTok{:}\NormalTok{n}
\NormalTok{y }\OtherTok{\textless{}{-}} \DecValTok{10}\SpecialCharTok{*}\NormalTok{x }\SpecialCharTok{+} \FunctionTok{rnorm}\NormalTok{(n, }\DecValTok{0}\NormalTok{, }\DecValTok{10}\NormalTok{)}
\NormalTok{n}
\end{Highlighting}
\end{Shaded}

\begin{verbatim}
[1] 30
\end{verbatim}

We can automatically say that the maximum value of the data is
294.6009735, or round it to a whole number: the maximum value of the
data is 295.

\subsection{Show some of the data}\label{show-some-of-the-data}

Let's combine the data in a tibble (think of it as a data frame if you
don't know what that is):

\begin{Shaded}
\begin{Highlighting}[]
\NormalTok{data }\OtherTok{\textless{}{-}}\NormalTok{ tibble}\SpecialCharTok{::}\FunctionTok{tibble}\NormalTok{(x, y)}
\NormalTok{data}
\end{Highlighting}
\end{Shaded}

\begin{verbatim}
# A tibble: 30 x 2
       x     y
   <int> <dbl>
 1     1  23.7
 2     2  14.4
 3     3  33.6
 4     4  46.3
 5     5  54.0
 6     6  58.9
 7     7  85.1
 8     8  79.1
 9     9 110. 
10    10  99.4
# i 20 more rows
\end{verbatim}

(only the first 10 rows get printed here thanks to it being a tibble).

To have a basic table, we can do

\begin{Shaded}
\begin{Highlighting}[]
\FunctionTok{kable}\NormalTok{(data[}\DecValTok{1}\SpecialCharTok{:}\DecValTok{10}\NormalTok{,])}
\end{Highlighting}
\end{Shaded}

\begin{longtable}[]{@{}rr@{}}
\toprule\noalign{}
x & y \\
\midrule\noalign{}
\endhead
\bottomrule\noalign{}
\endlastfoot
1 & 23.70958 \\
2 & 14.35302 \\
3 & 33.63128 \\
4 & 46.32863 \\
5 & 54.04268 \\
6 & 58.93875 \\
7 & 85.11522 \\
8 & 79.05341 \\
9 & 110.18424 \\
10 & 99.37286 \\
\end{longtable}

To make it look a bit better we can add an option

\begin{Shaded}
\begin{Highlighting}[]
\FunctionTok{kable}\NormalTok{(data[}\DecValTok{1}\SpecialCharTok{:}\DecValTok{10}\NormalTok{,]) }\SpecialCharTok{\%\textgreater{}\%}
  \FunctionTok{kable\_material}\NormalTok{(}\StringTok{"striped"}\NormalTok{)}
\end{Highlighting}
\end{Shaded}

\begin{longtable}[t]{rr}
\toprule
x & y\\
\midrule
1 & 23.70958\\
2 & 14.35302\\
3 & 33.63128\\
4 & 46.32863\\
5 & 54.04268\\
\addlinespace
6 & 58.93875\\
7 & 85.11522\\
8 & 79.05341\\
9 & 110.18424\\
10 & 99.37286\\
\bottomrule
\end{longtable}

You can set all sorts of formatting (and do it once at the start for all
tables) with the \texttt{kable} and \texttt{kableExtra} packages (see
the R Markdown book), so don't worry about formatting for now.

\subsection{Plot then fit a regression
data}\label{plot-then-fit-a-regression-data}

Now to plot the data:

\begin{Shaded}
\begin{Highlighting}[]
\FunctionTok{plot}\NormalTok{(x, y)}
\end{Highlighting}
\end{Shaded}

\pandocbounded{\includegraphics[keepaspectratio]{10.RMarkdownA_files/figure-latex/plot-1.pdf}}

To fit and then print the summary regression output from R:

\begin{Shaded}
\begin{Highlighting}[]
\NormalTok{fit }\OtherTok{\textless{}{-}} \FunctionTok{lm}\NormalTok{(y }\SpecialCharTok{\textasciitilde{}}\NormalTok{ x)}
\FunctionTok{print}\NormalTok{(}\FunctionTok{summary}\NormalTok{(fit))}
\end{Highlighting}
\end{Shaded}

\begin{verbatim}

Call:
lm(formula = y ~ x)

Residuals:
    Min      1Q  Median      3Q     Max 
-26.246  -4.503  -1.061   8.804  22.083 

Coefficients:
            Estimate Std. Error t value Pr(>|t|)    
(Intercept)   6.9136     4.5890   1.507    0.143    
x             9.5982     0.2585  37.131   <2e-16 ***
---
Signif. codes:  0 '***' 0.001 '**' 0.01 '*' 0.05 '.' 0.1 ' ' 1

Residual standard error: 12.25 on 28 degrees of freedom
Multiple R-squared:  0.9801,    Adjusted R-squared:  0.9794 
F-statistic:  1379 on 1 and 28 DF,  p-value: < 2.2e-16
\end{verbatim}

And for a report we can produce a simple table (including a caption) of
output and the regression fit:

\begin{Shaded}
\begin{Highlighting}[]
\FunctionTok{kable}\NormalTok{(}\FunctionTok{coefficients}\NormalTok{(}\FunctionTok{summary}\NormalTok{(fit)),}
      \AttributeTok{caption =} \StringTok{"Linear regression fit."}\NormalTok{) }\SpecialCharTok{\%\textgreater{}\%}
  \FunctionTok{kable\_material}\NormalTok{(}\StringTok{"striped"}\NormalTok{)}
\end{Highlighting}
\end{Shaded}

\begin{longtable}[t]{lrrrr}
\caption{\label{tab:fit}Linear regression fit.}\\
\toprule
 & Estimate & Std. Error & t value & Pr(>\&\#124;t\&\#124;)\\
\midrule
(Intercept) & 6.913647 & 4.5890369 & 1.506557 & 0.1431241\\
x & 9.598208 & 0.2584946 & 37.131179 & 0.0000000\\
\bottomrule
\end{longtable}

And create a plot:

\begin{Shaded}
\begin{Highlighting}[]
\FunctionTok{plot}\NormalTok{(x, y)}
\FunctionTok{abline}\NormalTok{(fit, }\AttributeTok{col=}\StringTok{"red"}\NormalTok{)}
\end{Highlighting}
\end{Shaded}

\pandocbounded{\includegraphics[keepaspectratio]{10.RMarkdownA_files/figure-latex/plotfit-1.pdf}}

\subsection{Now, go back and change the
data}\label{now-go-back-and-change-the-data}

The \emph{big} feature of dynamically generating reports is when you go
back and change or update the input data. For example, changing the data
in the above example and then re-running it to redo the report. The best
way to demonstrate this is for you to do it in the Exercise.

See the main document for what to do with this file for the Exercise.

\subsection{Automating text}\label{automating-text}

Once you've done the Exercise, you can even get a bit clever with your
writing by including an R \texttt{ifelse} statement to somewhat automate
the text, for example:

So the maximum value of \(y\) is 295, which is less than 400.

But you have to be careful and think about all possibilities -- what if
\(y=399.9\)? ```

Note that the \texttt{echo\ =\ FALSE} parameter was added to the code
chunk to prevent printing of the R code that generated the plot.

\end{document}
